% glossary: "Words we use" box for Sec. II (added 2026-10-07). One plain line per entry.
% The inline first-use explanations in the abstract, the Introduction and Sec. II are still the primary definitions.
\newcommand{\gw}[2]{\textbf{#1}: #2\par\vspace{0.6pt}}
\onecolumngrid
\noindent\begin{minipage}{\textwidth}
\vspace{4pt}\par\noindent\refstepcounter{table}\label{tab:words}{\small TABLE~\thetable. Words we use, in plain terms.}\par\vspace{3pt}
\noindent\fcolorbox{nullc}{white}{%
\begin{minipage}[t]{0.465\textwidth}\scriptsize\raggedright
\gw{Agent}{one language model inside a scaffold; it uses a computer and a group chat.}
\gw{Scaffold}{the program that runs an agent: it builds the model's input and carries out its output.}
\gw{Model call (call)}{one turn: the model reads its context and writes one output (a message, an action or a wait).}
\gw{Context window}{the text that a call reads; the agent's short-term memory.}
\gw{Forced erasure (wipe, reset)}{the scaffold empties the context window at a fixed size (41 records in regime III); notes and files survive.}
\gw{Call clock}{time counted in an agent's own calls, not in minutes. Ten idle minutes between two calls add nothing.}
\gw{Read-out call}{the first call whose context holds a given message: the moment the agent sees it.}
\gw{In flight}{posted, but not yet in the reader's context. Our placebo.}
\gw{Named message}{a message that mentions the reader by name.}
\gw{Goal period (\#1--\#51)}{the time for one shared goal, usually a week.}
\gw{Kickoff}{the message that announces a goal.}
\gw{Room}{a separate chat channel.}
\gw{Regime I, II, III}{the three eras of the scaffold: one room in daily sessions; several rooms in sessions; continuous, own room only, forced erasures.}
\gw{Scheduler}{the operator's daily timetable that turns agents on and off.}
\gw{Natural experiment (NE\#\#)}{a dated step change in the setup (scaffold, roster, rooms, operator), tested before vs.\ after; not a goal period, and often spans two.}
\end{minipage}}\hfill
\fcolorbox{nullc}{white}{%
\begin{minipage}[t]{0.465\textwidth}\scriptsize\raggedright
\gw{Nudge}{an automatic prompt to an idle agent.}
\gw{Unit}{a goal period, or the part of one between two scaffold changes.}
\gw{Impostor}{a shared input that fakes influence between agents.}
\gw{Null}{the simplest model with no effect; a claim must beat it.}
\gw{Spin}{a part with a few states; here an agent that talks, works or waits.}
\gw{Field}{an input that acts on each agent directly: timetable, goal, room, style.}
\gw{Coupling}{the effect of one agent on another, through a message that it reads.}
\gw{Quench}{a sudden change of the field, such as a kickoff.}
\gw{Loop gain $g$}{new messages that one message causes directly. Below 1, echoes die out.}
\gw{Fano factor}{variance of a count over its mean: 1 for purely random events, more for bursts.}
\gw{Collective mode}{a pattern in which many agents move together.}
\gw{Marchenko--Pastur edge}{the largest correlation that chance alone gives in a record of this size.}
\gw{Well / kick}{an agent's slow resting point / the short push of one message read.}
\gw{Reserved data}{goal periods set aside before exploration, for confirmation tests.}
\gw{M02, \ldots, M17}{a physics model, named by its folder number (Table~\ref{tab:models}).}
\gw{H01, \ldots, H142}{a hypothesis card: one pre-registered test, with its data, prediction and score (Appendix~\ref{app:eu}). Set small and gray in the text.}
\gw{Credence}{Claude's judged probability that a claim would survive a test on reserved data; not a frequency.}
\gw{bge, gte}{two text-embedding models; we report content results with both.}
\gw{Trim}{keep only the hours when every agent is switched on, so the timetable cannot fake agreement.}
\gw{Post hoc}{found after the test, not predicted; reported, but not counted as evidence.}
\end{minipage}}
\end{minipage}\par\vspace{6pt}
\twocolumngrid
